There has been much recent progress in sequence modeling through recurrent neural networks (RNN; \citealt{sutskever2014sequence,bahdanau2015neural,wu2016google}), convolutional networks (CNN; \citealt{kalchbrenner2016nmt,gehring2016convolutional,gehring2017convs2s,kaiser2017depthwise}) and self-attention models \citep{paulus17intra,vaswani2017transformer}. 
RNNs integrate context information by updating a hidden state at every time-step, CNNs summarize a fixed size context through multiple layers, while as self-attention directly summarizes all context.

Attention assigns context elements \emph{attention weights} which define a weighted sum over context representations \citep{bahdanau2015neural,sukhbaatar2015memnet,chorowski2015asr,luong2015effective}.
Source-target attention summarizes information from another sequence such as in machine translation while as self-attention operates over the current sequence.
Self-attention has been formulated as \emph{content-based} where attention weights are computed by comparing the current time-step to all elements in the context (Figure~\ref{fig:sattn2}). 
The ability to compute comparisons over such unrestricted context sizes are seen as a key characteristic of self-attention \citep{vaswani2017transformer}.

\begin{figure}[h]
\centering
\begin{subfigure}[b]{0.45\textwidth}
\centering
\includegraphics[width=1\textwidth]{fig/content.png}
\caption{Self-attention}
\label{fig:sattn2}
\end{subfigure}
\qquad \quad
\begin{subfigure}[b]{0.45\textwidth}
\centering
\includegraphics[width=1\textwidth]{fig/location.png}
\caption{Dynamic convolution}
\label{fig:dsdconv}
\end{subfigure}
\caption{Self-attention computes attention weights by comparing all pairs of elements to each other (\subref{fig:sattn2}) while as dynamic convolutions predict separate kernels for each time-step (\subref{fig:dsdconv}).}
\label{fig:attention}
\end{figure}

However, the ability of self-attention to model long-range dependencies has recently come into question \citep{tang2018why} and the unlimited context size is computationally very challenging due to the quadratic complexity in the input length. 
Furthermore, in practice long sequences require the introduction of hierarchies \citep{liu2018wikisum}. 

In this paper, we introduce \emph{\sdconvfull{}s} which are depth-wise separable \citep{sifre2014rigid,Chollet2017XceptionDL,kaiser2017depthwise}, softmax-normalized and share weights over the channel dimension.
The result is a convolution with several orders of magnitude fewer weights than a standard non-separable convolution. 
Different to self-attention, \sdconvfull{}s reuse the same weights for context elements, regardless of the current time-step.

\emph{\Tvsdconvfull{}s} build on \sdconvfull{}s by predicting a different convolution kernel at every time-step. The kernel is a function of the current time-step only as opposed to the entire context as in self-attention (Figure~\ref{fig:dsdconv}).
\Tvsdconvfull{}s are similar to locally connected layers in the sense that the weights change at every position, however, the difference  is that weights are dynamically generated by the model rather than fixed after training \citep{lecun1998grad,taigman2014deepface,chen2015interspeech}.
Our approach also bears similarity to location-based attention which does not access the context to determine attention weights, however, we do not directly take the attention weights from the previous time-step into account \citep{chorowski2015asr,luong2015effective}.
\citet{shen2018block} reduce complexity by performing attention within blocks of the input sequence and 
\citet{shen2017disan,shen2018fast} perform more fine-grained attention over each feature.
\citet{shen2018emnlp} and \citet{gong2018emnlp} use input-dependent filters for text classification tasks.


Our experiments show that \sdconvfull{}s perform competitively to strong self-attention results and that \tvsdconvfull{}s can perform even better.
On WMT English-German translation \tvsdconvfull{}s achieve a new state of the art of 29.7 BLEU, on WMT English-French they match the best reported result in the literature, and on IWSLT German-English \tvsdconvfull{}s outperform self-attention by 0.8 BLEU.
\Tvsdconvfull{}s achieve 20\% faster runtime than a highly-optimized self-attention baseline.
For language modeling on the \gbw{} benchmark \tvsdconvfull{}s perform as well as or better than self-attention and on \cnndm{} abstractive document summarization we outperform a strong self-attention model.


